\newpage \begin{bclogo}[barre=none,logo={}]{Contraction du texte} 
	\begin{center}
		\subsubsection*{LE CERVEAU A INVENTER LA MACHINE \emph{Comment naissent les grandes découvertes} $$\sim\ \text{{\bfseries Yves AGID}}\ \sim$$ }\addcontentsline{toc}{section}{
				20240712 - Contraction de texte (AS 2024)
			}
	\end{center}
\end{bclogo}
	
	\hspace*{0.5cm}L'homme est avide de progrès \textendash\ revêtant un caractère multidimensionnel \textendash\ car il est déterminant pour sa survie. Le progrès est en grande partie la résultante d'un savoir basé sur l'existant. En outre, l'aptitude naturelle de l'humain à comprendre et à créer sont pour lui des outils nécessaires. Cette avidité de progrès découle aussi de la soif humaine pour la connaissance de son environnement. \newline
	
	\hspace*{0.5cm}La recherche est laborieuse et décourageante pour tout scientifique. C'est le propre même de la science. Ainsi, les meilleurs scientifiques sont les moins défaitistes et les plus déterminés à prouver leurs prédicats. Un exemple notable est la découverte des hiéroglyphes par Jean-François Champollion et du vaccin ARN messager contre le virus Covid-19 par Katalin Karikö respectivement après quatorze et trente années de recherche.\newline
	
	\hspace*{0.5cm}La découverte nécessite de percevoir l'invisible. D'où la mise au point d'instruments sophistiqués (microscopes, télescopes...) capables d'identifier l'imperceptible. Aussi, l'ordinateur et l’intelligence artificielle aident l'\emph{Homo sapiens} à comprendre le phénomène identifié bien qu'ils ne soient que des représentations simplistes du cerveau humain. Cependant, cela n'empêche pourtant pas le mathématicien Étienne Ghys de faire une analogie entre Internet et le cerveau humain.\newline
	
	\hspace*{0.5cm}Le foisonnement de l'information scientifique collectivise la science et révèle une catégorie de scientifiques sachant détecter l'information pertinente à même de déboucher sur une découverte, parmi les nombreuses idées scientifiques, découlant souvent d'équipes de recherche nécessitant inéluctablement un leader pour éviter toutes dérives. Toutefois, la proportion des découvertes d'une région dépend du contexte socio-économique. Néanmoins, tout bon scientifique se devra donc d'\oe uvrer avec la communauté scientifique en apportant des idées nouvelles pour l'avancée scientifique.
	
\begin{center}
	\textbf{282 mots}
\end{center}$$\sim\star\star\star\sim$$

